\section{Discussion}

% Plan: State the methodological implication and identify the paper's unit of analysis.
The central implication is methodological: answer accuracy is not an identifying
measurement of belief use. A model can preserve accuracy while changing which
described source controls the answer, or lose accuracy because a linguistic
construction restricts access to the correct source. Typed competitors turn
that ambiguity into an estimable interaction. For linguistically grounded
evaluation, the relevant unit is therefore not a paraphrase average but a
triple: the intended source, the competing source types, and the construction
that regulates access between description and answer.

% Plan: Interpret the held-out residual as a composition result while preserving the operator boundary.
The second implication is theoretical. Linguistic restrictions need not compose
as independent penalties. In Qwen2.5-7B-Instruct, voice and embedding constraints
share a behavioral bottleneck: their joint cost is smaller than the sum of their
isolated costs. This is stronger than documenting another sensitivity because
the additive account predicts unseen compounds and fails quantitatively. At the
same time, positive residuals identify only a subadditive class. The data do not
distinguish a weakest-link rule from every smooth saturating alternative, and
they do not imply that the model represents accessibility as a scalar variable.

% Plan: Explain the checkpoint result as a boundary on transfer rather than a causal training claim.
The checkpoint comparison adds a third implication. The typed source-control
object remains useful outside the positive checkpoint, but its composition law
is not portable in direction or magnitude. Because affine readout scaling would
preserve residual sign and normalized geometry, the failed transfer indicates a
functional change in how linguistic constraints reach the answer interface.
The matched Base contrast makes instruction tuning a plausible source of that
change, but only an association rather than a causal training result.

% Plan: State limitations as claim boundaries, not as a defensive experiment log.
These boundaries define the paper's scope. The stimuli are controlled synthetic
English worlds, the answer interface uses audited single-token candidates, and
the dependent variable is model behavior rather than human interpretation. We
therefore do not claim unrestricted English generality, broad model-family
transfer, human source-monitoring equivalence, or a hidden-state belief circuit.
The linguistic validity audit further treats grammaticality, source identity,
commitment, and target invariance as properties to be independently checked,
not assumed from generator intent.

% Plan: Close with the single transferable research question and the next scientific test.
\section{Conclusion}

This paper reframes linguistic sensitivity in theory-of-mind evaluation as a
source-binding problem. Typed conflict probes make the answer controller
identifiable, a frozen compound test reveals subadditive composition, and the
checkpoint comparison shows that this law is a boundary-sensitive behavioral
property rather than a portable model-family invariant. The next test is not
another paraphrase grid, but a controlled transfer of the typed geometry across
models or languages while freezing the same primitive costs. The central
question is therefore which linguistic systems make a source govern an answer,
and how those systems compose.
